\subsection{The natural right-resolving presentation}

We now return to the subset construction.  The exact decoding theorem shows
that the determinized presentation has a finite transient ambiguity shell and
that its essential core is already visible in the explicit graph $G_m$.  The
next results quantify the resulting deterministic presentation size on the
natural folded alphabet $X_m$.

\begin{definition}[Deterministic candidate-set presentation]
\label{def:Phi-m-right-resolving}
For a nonempty subset $S\subset V_m$ and a label $a\in X_m$, define
\[
\delta_m(S,a):=
\{v'\in V_m:\exists\, e:v\to v' \text{ with } v\in S \text{ and } \ell_m(e)=a\}.
\]
Starting from the state $V_m$, take all reachable nonempty subsets of $V_m$ and
add labeled edges
\[
S\xrightarrow{a}\delta_m(S,a)
\]
whenever $\delta_m(S,a)\neq \varnothing$.  The resulting finite labeled graph
is denoted $H_m^{\mathrm{det}}$.
\end{definition}

\begin{proposition}[Deterministic right-resolving presentation]
\label{prop:Phi-m-right-resolving}
The graph $H_m^{\mathrm{det}}$ is right-resolving and presents the same labeled
language as $G_m$.  Consequently its edge shift has label image $Y_m$.
\end{proposition}

\begin{proof}
By construction, for a fixed state $S$ and label $a$, there is at most one
outgoing edge with label $a$, so $H_m^{\mathrm{det}}$ is right-resolving.  The
subset construction preserves the accepted labeled language of a finite
automaton, hence $H_m^{\mathrm{det}}$ and $G_m$ present the same set of finite
label words.  Their bi-infinite label shifts therefore coincide, and by
Theorem~\ref{thm:Phi-m-sofic-graph} this common image is $Y_m$.
\end{proof}

\begin{proposition}[The singleton core]
\label{prop:deterministic-singleton-core}
Let $m\ge 3$.  Then:
\begin{enumerate}[label=\textup{(\roman*)}]
\item every state of $H_m^{\mathrm{det}}$ reached after at least $m$ input
symbols is a singleton;
\item the reachable singleton states are exactly the sets $\{v\}$ with
$v\in V_m$;
\item the labeled subgraph induced by the singleton states is canonically
labeled-isomorphic to $G_m$.
\end{enumerate}
\end{proposition}

\begin{proof}
Take any block
\[
a_1\cdots a_m\in\mathcal L_m(Y_m).
\]
By Theorem~\ref{thm:length-m-block-unique-lift}, this block has a unique raw
lift, hence a unique length-$m$ path in $G_m$ carrying that label word.  Its
terminal vertex is therefore uniquely determined by the block.  Consequently,
every reachable state after reading $m$ symbols is a singleton.  Since
singletons are forward closed under the subset construction, every state reached
after at least $m$ symbols is again a singleton.  This proves
\textup{(i)}.

For \textup{(ii)}, it remains to show that every singleton $\{v\}$ with
$v\in V_m$ is reachable.  By
Theorem~\ref{thm:alternating-single-symbol-synchronizing}, the symbol $b_m$
sends the full candidate set to a singleton.  The
underlying directed graph of $G_m$ is the binary de Bruijn graph on words of
length $m-1$, so it is strongly connected.  Thus for every $v\in V_m$ there
exists a path in $G_m$ from that singleton state to $v$, and reading its label
from the singleton yields $\{v\}$.

Finally, \textup{(iii)} is immediate from the definitions: for singleton states,
\[
\delta_m(\{v\},a)=\{v'\}
\]
if and only if $G_m$ has an edge $v\to v'$ labeled $a$.  Hence the singleton
subgraph is exactly a relabeling of $G_m$.
\end{proof}

\begin{proposition}[The ambiguity shell is transient]
\label{prop:deterministic-ambiguity-transient}
Let $m\ge 3$.  Every non-singleton state of $H_m^{\mathrm{det}}$ is transient.
Equivalently, the subgraph induced by the non-singleton states is acyclic.
\end{proposition}

\begin{proof}
Suppose there were a directed cycle consisting entirely of non-singleton states.
Since every state of $H_m^{\mathrm{det}}$ is reachable from the initial state
$V_m$, one could first reach that cycle and then traverse it often enough to
obtain a path of length at least $m$ ending at a non-singleton state.  This
contradicts Proposition~\ref{prop:deterministic-singleton-core}\textup{(i)}.
Hence no such cycle exists, and every non-singleton state is transient.
\end{proof}

\begin{proposition}[Exact nilpotency of the ambiguity shell]
\label{prop:ambiguity-shell-nilpotency}
Let $m\ge 3$.  Order the states of $H_m^{\mathrm{det}}$ by placing the
non-singleton states first and the singleton core last.  Then the adjacency
matrix of $H_m^{\mathrm{det}}$ has the form
\[
A(H_m^{\mathrm{det}})
=
\begin{pmatrix}
N_m & R_m\\
0 & A(G_m)
\end{pmatrix}
\]
for some matrix $R_m$, and
\[
N_m^m=0,
\qquad
N_m^{m-1}\neq 0.
\]
Consequently
\[
\det(I-zA(H_m^{\mathrm{det}}))=\det(I-zA(G_m)).
\]
\end{proposition}

\begin{proof}
By Proposition~\ref{prop:deterministic-singleton-core}\textup{(iii)}, the
singleton states induce a labeled copy of $G_m$.  Since singleton states are
forward closed, there are no edges from the singleton core back into the
non-singleton shell, so the adjacency matrix has the displayed block upper
triangular form.

By Proposition~\ref{prop:deterministic-singleton-core}\textup{(i)}, every path
of length $m$ in $H_m^{\mathrm{det}}$ ends in a singleton.  Hence there is no
path of length $m$ contained entirely in the non-singleton shell, which is
equivalent to $N_m^m=0$.

By Corollary~\ref{cor:uniform-synchronizing-length}, there exists a
non-synchronizing block of length $m-1$.  Reading that block from the initial
state $V_m$ yields a path of length $m-1$ ending at a non-singleton state.
Because singleton states are forward closed, every intermediate state on that
path is also non-singleton.  Therefore the shell contains a path of length
$m-1$, so $N_m^{m-1}\neq 0$.

Finally,
\[
\det(I-zA(H_m^{\mathrm{det}}))
=
\det(I-zN_m)\det(I-zA(G_m)).
\]
Since $N_m$ is nilpotent, $\det(I-zN_m)=1$, and the determinant identity
follows.
\end{proof}

\begin{proposition}[The explicit graph is follower separated]
\label{prop:G-m-follower-separated}
Let $m\ge 3$.  Distinct vertices of $G_m$ have distinct follower sets.
\end{proposition}

\begin{proof}
Let
\[
b_m=u_1\cdots u_m
\]
be the alternating symbol from
Definition~\ref{def:alternating-extremal-words}.  For each vertex
\[
v=v_1\cdots v_{m-1}\in V_m,
\]
consider the length-$m$ path obtained by appending the fixed bit word
$u_1\cdots u_m$.  Let
\[
\Lambda(v)\in\mathcal L_m(Y_m)
\]
denote its label word.

We claim that $\Lambda(v)$ is carried by exactly one path in $G_m$, namely the
one that starts at $v$ and appends $u_1,\ldots,u_m$.  Indeed, the last symbol
of $\Lambda(v)$ is $b_m$, and by
Theorem~\ref{thm:alternating-single-symbol-synchronizing}, the symbol $b_m$ is
carried by a unique edge.  Thus the final edge of any path carrying
$\Lambda(v)$ is fixed.  Since $G_m$ is one-step left-resolving by
Proposition~\ref{prop:one-step-left-resolving}, the penultimate label then
determines the penultimate edge uniquely; iterating this backward reconstruction
shows that the whole path carrying $\Lambda(v)$ is unique.

In particular, its initial vertex is uniquely determined by $\Lambda(v)$.
Hence if $v\neq w$, the word $\Lambda(v)$ belongs to the follower set of $v$
but not to that of $w$.  Therefore distinct vertices have distinct follower
sets.
\end{proof}

\begin{corollary}[The graph $G_m$ is the canonical minimal right-resolving presentation]
\label{cor:G-m-fischer-cover}
For every $m\ge 3$, the canonical minimal right-resolving presentation of
$Y_m$ is labeled-isomorphic to $G_m$.  Equivalently, $G_m$ is the right
Fischer cover of $Y_m$.
\end{corollary}

\begin{proof}
By Theorem~\ref{thm:Phi-m-conjugacy-threshold}, the shift $Y_m$ is irreducible
(indeed conjugate to the full two-shift).  Proposition~\ref{prop:Phi-m-right-resolving}
gives a deterministic right-resolving presentation of $Y_m$.  By
Proposition~\ref{prop:deterministic-ambiguity-transient}, the only essential
component of that presentation is the singleton core from
Proposition~\ref{prop:deterministic-singleton-core}, which is labeled-isomorphic
to $G_m$.  Proposition~\ref{prop:G-m-follower-separated} shows that this core is
already follower separated.  The standard construction of the canonical
minimal right-resolving presentation, equivalently of the right Fischer cover,
as the follower-separated essential component of a deterministic
right-resolving presentation therefore leaves $G_m$ unchanged; see
\cite[Chs.~3--4]{LindMarcus1995} and \cite[Ch.~4]{Kitchens1998SymbolicDynamics}.
\end{proof}

\begin{theorem}[Size of the right-resolving presentation]
\label{thm:natural-state-complexity-gap}
For every $m\ge 3$, the canonical minimal right-resolving presentation of
$Y_m$ over the folded alphabet $X_m$ has exactly
\[
\lvert V(G_m)\rvert=2^{m-1},
\qquad
\lvert E(G_m)\rvert=2^m
\]
states and edges.  In particular, no right-resolving presentation of $Y_m$
over $X_m$ uses fewer than $2^{m-1}$ states.  Since
\[
\lvert X_m\rvert=F_{m+2},
\]
one has
\[
\frac{\lvert V(G_m)\rvert}{\lvert X_m\rvert}
=
\frac{2^{m-1}}{F_{m+2}}
\sim
\frac{\sqrt5}{2\varphi^2}\Bigl(\frac{2}{\varphi}\Bigr)^m,
\qquad
\varphi=\frac{1+\sqrt5}{2}.
\]
Thus, although every $Y_m$ with $m\ge 3$ is conjugate to the full two-shift,
its canonical deterministic presentation over the natural folded alphabet has
exponentially growing state complexity, even relative to the alphabet size.
\end{theorem}

\begin{proof}
Corollary~\ref{cor:G-m-fischer-cover} identifies $G_m$ as the canonical minimal
right-resolving presentation of $Y_m$.  The state and edge counts are therefore
those of $G_m$, and Definition~\ref{def:G-m} gives
\[
\lvert V(G_m)\rvert=2^{m-1},
\qquad
\lvert E(G_m)\rvert=2^m.
\]
Minimality of the Fischer cover yields the state lower bound for all
right-resolving presentations over $X_m$.

The alphabet size is
\[
\lvert X_m\rvert=F_{m+2}
\]
by Corollary~\ref{cor:stable-syntax-counting}.  Using the Binet asymptotic
$F_{m+2}\sim \varphi^{m+2}/\sqrt5$ gives
\[
\frac{2^{m-1}}{F_{m+2}}
\sim
\frac{2^{m-1}\sqrt5}{\varphi^{m+2}}
=
\frac{\sqrt5}{2\varphi^2}\Bigl(\frac{2}{\varphi}\Bigr)^m,
\]
which proves the claimed representation gap.
\end{proof}

\begin{corollary}[Coincidence of decoding and presentation scales]
\label{cor:local-complexity-scale-coincidence}
For every $m\ge 3$, the following quantities are exact and all equal to the
same scale:
\begin{enumerate}[label=\textup{(\roman*)}]
\item the finite-type order of $Y_m$ over the alphabet $X_m$ is $m$;
\item the minimal synchronizing length of the presentation $G_m$ is $m$;
\item the nilpotency index of the ambiguity-shell matrix $N_m$ is $m$;
\item the minimal memory of a one-sided sliding-block inverse of $\Phi_m$ is
$m-1$.
\end{enumerate}
In particular, although the shifts $Y_m$ with $m\ge 3$ are all conjugate to the
full two-shift, their natural presentations over the folded alphabets $X_m$
still reflect the window length.  Theorem~\ref{thm:exact-markov-order-bernoulli}
shows that the same threshold is also the Markov memory under every transported
Bernoulli law, while
Theorem~\ref{thm:natural-state-complexity-gap} quantifies the accompanying
exponential presentation gap on the natural alphabet.
\end{corollary}

\begin{proof}
Part \textup{(i)} is Theorem~\ref{thm:Y-m-exact-SFT-order},
part \textup{(ii)} is Corollary~\ref{cor:uniform-synchronizing-length},
part \textup{(iii)} is Proposition~\ref{prop:ambiguity-shell-nilpotency}, and
part \textup{(iv)} is Corollary~\ref{cor:Phi-m-memory-threshold-exact}.
\end{proof}
